
\documentclass{article}
\usepackage{colortbl}
\usepackage{makecell}
\usepackage{multirow}
\usepackage{supertabular}

\begin{document}

\newcounter{utterance}

\centering \large Interaction Transcript for game `ifeval', experiment `ifeval', episode 295 with Bentel rockers\_\_qwen3.5{-}4b\_\_Bentel\_iter\_3.
\vspace{24pt}

{ \footnotesize  \setcounter{utterance}{1}
\setlength{\tabcolsep}{0pt}
\begin{supertabular}{c@{$\;$}|p{.15\linewidth}@{}p{.15\linewidth}p{.15\linewidth}p{.15\linewidth}p{.15\linewidth}p{.15\linewidth}}
    \# & \multicolumn{2}{c}{Player} && \multicolumn{2}{c}{Game Master} \\
    \hline

    \theutterance \stepcounter{utterance}  
    & & & \multicolumn{4}{p{0.6\linewidth}}{
        \cellcolor[rgb]{0.9,0.9,0.9}{
            \makecell[{{p{\linewidth}}}]{
                \texttt{\tiny{[P1$\langle$GM]}}
                \texttt{Translate the following sentence into German and then criticize it: Werner was a good friend of mine, but not very smart.} \\
\texttt{Avoid the word "schlau" throughout your response.} \\
            }
        }
    }
     \\ \\

    \theutterance \stepcounter{utterance}  
    & \multicolumn{4}{p{0.6\linewidth}}{
        \cellcolor[rgb]{0.561,0.737,0.561}{
            \makecell[{{p{\linewidth}}}]{
                \texttt{\tiny{[P1$\rangle$GM]}}
                \texttt{**Übersetzung:**} \\
\texttt{Werner war ein guter Freund von mir, aber nicht sehr klug.} \\
\\ 
\texttt{**Kritik:**} \\
\texttt{Die Übersetzung ist grammatikalisch korrekt und die Bedeutung wird präzise übertragen. Allerdings ist die Wortwahl bei "nicht sehr klug" etwas direkt und möglicherweise unangenehm, da sie eine negative Eigenschaft des Freundes sofort benennt. In einer deutschen Gesprächssituation könnte man dies vielleicht etwas abmildern, indem man den Fokus auf andere Aspekte seiner Persönlichkeit legt oder den Satz so formuliert, dass die Unzulänglichkeit weniger persönlich wirkt, ohne dabei den Kern der Aussage zu verlieren. Zudem ist die Struktur "guter Freund ... aber nicht sehr klug" ein klassisches Beispiel für einen Kontrast, der in der deutschen Sprache sehr natürlich klingt, solange der Kontext dies erlaubt.} \\
            }
        }
    }
    & & \\ \\

    \theutterance \stepcounter{utterance}  
    & & & \multicolumn{2}{p{0.3\linewidth}}{
        \cellcolor[rgb]{0.9,0.9,0.9}{
            \makecell[{{p{\linewidth}}}]{
                \texttt{\tiny{[GM$|$GM]}}
                {'keywords:forbidden_words': {'forbidden_words': ['schlau']}}
            }
        }
    }
    & & \\ \\

    \theutterance \stepcounter{utterance}  
    & & & \multicolumn{2}{p{0.3\linewidth}}{
        \cellcolor[rgb]{0.9,0.9,0.9}{
            \makecell[{{p{\linewidth}}}]{
                \texttt{\tiny{[GM$|$GM]}}
                \texttt{game\_result = WIN} \\
            }
        }
    }
    & & \\ \\

\end{supertabular}
}

\end{document}
